\subsection{The runs that were summarised}

Two red-oak runs and the slope run were described in prose so that the main tables would stay readable. They are recorded here at the same precision as Table~\ref{tab:redmild}, because a summary that cannot be checked against the script is not a result. Table~\ref{tab:rest} collects them. The both-mode frequency is 9~Hz and the vacuum is 0 in the slow and slope runs, and 3~Hz and 0 in the harsh run already discussed.

\begin{table}[htbp]
\centering
\caption{Remaining red-oak runs. Speed is in m/s, slope in degrees, terrain frequency in Hz.}
\label{tab:rest}
\begin{tabular}{@{}llrrrrr@{}}
\toprule
Case & Mode & $Q_{\mathrm{f}}$ & $I_{\mathrm{mult}}$ & $I_{\mathrm{miss}}$ & $C_{\mathrm{v}}$ & $\eta_{\mathrm{E}}$ \\
\midrule
0.70~m/s, 0$^\circ$, 6~Hz & Ground & 58.02 & 21.53 & 20.45 & 19.45 & 51.45 \\
 & PID & 60.24 & 21.35 & 18.41 & 21.44 & 59.48 \\
 & Vision & 65.58 & 13.27 & 21.16 & 18.12 & 57.41 \\
 & MPC & 77.49 & 12.09 & 10.42 & 16.11 & 77.68 \\
 & Both & 88.17 & 5.84 & 5.99 & 13.09 & 87.31 \\
\midrule
1.10~m/s, 12$^\circ$, 6~Hz & Ground & 46.34 & 20.87 & 32.79 & 17.57 & 28.89 \\
 & PID & 60.62 & 22.17 & 17.21 & 20.40 & 60.49 \\
 & Vision & 68.46 & 13.43 & 18.11 & 18.11 & 62.01 \\
 & MPC & 77.28 & 11.46 & 11.26 & 16.01 & 75.82 \\
 & Both & 88.22 & 5.61 & 6.17 & 12.57 & 87.05 \\
\bottomrule
\end{tabular}
\end{table}

The slow run repeats the mild-run pattern at a lower shaft speed. Vision lowers the multiple index from 21.4 to 13.3 percent and lowers efficiency from 59.5 to 57.4 percent. The rate side effect is visible even when the machine is already going slowly. MPC then both restore efficiency, to 77.7 and 87.3 percent. Nothing in the slow run suggests that the stack depends on the 1.10~m/s condition.

The slope run is the cleanest demonstration of requirement~2. Ground-drive efficiency is 28.9 percent with a miss index of 32.8 percent. Every electronic mode returns to within about one point of its level-ground efficiency at the same speed. The both-mode efficiency is 87.1 percent on the slope and 87.0 percent on the level. The model has no other slope channel. The table is what that modelling choice looks like when it is not smoothed into a sentence. A field trial in which electronic control failed to recover a slope would be evidence of a channel the model omitted, most likely opener draft, and not evidence that the arithmetic of Table~\ref{tab:rest} was mis-copied.

\subsection{Reading the near-tie under harsh vibration}

Section~\ref{sec:estimation} reported a 3~Hz command for red oak at 16~Hz terrain. The grid slice in Chapter~\ref{ch:design} shows why that command should not be over-interpreted. At the nominal diameter the predicted costs at 3~Hz and at 9~Hz differ by about $0.0001$. The drawn sample mean in the recorded run selected 3~Hz. Plant-side, with $\alpha_{\mathrm{iso}}=0.55$, a 3~Hz command and a 16~Hz terrain give $f_{\mathrm{eff}}\approx 9.3$~Hz, and a 9~Hz command would give about 12.6~Hz. The plant therefore experiences a real difference between the two commands even though the predictor, using the heavier weight 0.35 on $f_{\mathrm{dist}}^2$, scores them almost alike. The run that was printed used 3~Hz. A second run whose sample mean broke the tie the other way would print 9~Hz and would show a higher plant-side harshness. The efficiency gain over PID would remain, because both corners beat a command of 0~Hz, which is the PID and ground setting. The identity of the winning hertz would not remain. That is the correct confidence statement. The coordinator rejects the middle of the frequency range and rejects a non-zero vacuum. It does not, on this cost, identify a unique frequency once the terrain is already harsh.

\subsection{A note on repeatability of the other scenarios}

The twelve-repeat standard deviation of 0.89 points was computed only for the mixed-lot efficiency gain. The other scenarios were single runs of 4{,}000 events. A binomial picture of a proportion near 0.6 with 4{,}000 trials has a standard error near $\sqrt{0.6\cdot 0.4/4000}\approx 0.008$, or 0.8 percentage points, before the dependence among successive spacings is considered. Dependence widens the error. Doubling the binomial picture is a cautious reading: differences of two points in the tables are inside noise, and differences of ten points are not. The vision-versus-PID multiple-index drop on mild red oak is about ten points. The both-versus-ground efficiency gaps are tens of points. The vision-versus-PID efficiency differences of one point are not a finding. They were described as ``does not follow'' and ``stays'', which is the language that one point deserves.

The script seed is part of the result. A reader who changes the seed will move the fourth digit of a long run and can move the harsh-case frequency across the tie. Publishing the seed, 20261029, and the repeat rule $20261029+17r$, is what makes that movement a check rather than a contradiction. The file \texttt{estimation\_results.json} stores the run that the tables were typed from.

\subsection{Consolidated efficiency}

Table~\ref{tab:all} puts every scenario's planting efficiency on one page so that the ranking does not depend on which prose paragraph a reader happens to remember. The both-mode column is the stack. The ground column is the open-loop machine. The gap between them is the quantity the title of this manuscript refers to, and it is a model quantity.

\begin{table}[htbp]
\centering
\caption{Planting efficiency $\eta_{\mathrm{E}}$ in percent for every estimated scenario. Terrain frequency $f_{\mathrm{dist}}$ is in hertz and slope is in degrees. Speed is in m/s.}
\label{tab:all}
\begin{tabular}{@{}llrrrrr@{}}
\toprule
Lot & Condition & Ground & PID & Vision & MPC & Both \\
\midrule
Red oak & 0.70, 0$^\circ$, 6~Hz & 51.45 & 59.48 & 57.41 & 77.68 & 87.31 \\
Red oak & 1.10, 0$^\circ$, 6~Hz & 52.38 & 60.97 & 60.70 & 76.56 & 87.00 \\
Red oak & 1.10, 12$^\circ$, 6~Hz & 28.89 & 60.49 & 62.01 & 75.82 & 87.05 \\
Red oak & 1.10, 0$^\circ$, 16~Hz & 17.11 & 21.77 & 21.99 & 64.46 & 64.28 \\
Pin oak & 1.10, 0$^\circ$, 6~Hz & 54.87 & 61.61 & 63.78 & 75.20 & 86.49 \\
Bur oak & 0.70, 0$^\circ$, 6~Hz & 14.76 & 18.35 & 16.85 & 45.62 & 57.86 \\
Bur oak & 1.10, 8$^\circ$, 12~Hz & 5.01 & 11.93 & 9.82 & 43.99 & 55.12 \\
Mixed & 1.10, 8$^\circ$, 12~Hz & 20.12 & 37.75 & 33.52 & 65.19 & 72.25 \\
\bottomrule
\end{tabular}
\end{table}

The column order is the order in which a builder would add the pieces. PID recovers the slope row and little else. Vision, added to PID, does not dominate that column: it wins on pin oak and on the slope, and it loses a point or more on the slow red oak, on both bur-oak rows, and on the mixed lot. MPC, added instead of vision, wins every row against PID by a wide margin, and the margin is widest where the ground drive has collapsed. Both, added last, wins every row against MPC except the harsh red-oak row, where it is 0.18 points lower, inside the repeatability note above. The architecture's claim is not that the stack is uniformly better than its better half. The claim is that the better half is different in different rows, and that publishing both indices, of which this table shows only efficiency, keeps a builder from turning on a loop that the row does not need. On a harsh red-oak bout the factor loop is the one that pays. On a pin-oak bout the camera still earns its place, because the both-mode efficiency of 86.5 percent is eleven points above MPC. The table is the decision.

